% ============================================================================
% ANEXOS
% ============================================================================

\chapter*{Anexos}
\addcontentsline{toc}{chapter}{Anexos}

% ============================================================================
% ANEXO B: TABLA COMPLETA ESTADO DEL ARTE
% ============================================================================

\section{Análisis Comparativo Exhaustivo de Investigaciones Previas}
\label{sec:anexo_tabla_completa}

Esta secci├│n presenta el an├ílisis comparativo exhaustivo de 17 investigaciones representativas sobre clasificaci├│n de comportamiento sedentario y actividad f├¡sica mediante dispositivos wearables. La tabla abarca diversas t├®cnicas metodol├│gicas (l├│gica difusa, machine learning, deep learning, m├®todos estad├¡sticos), tipos de validaci├│n (destacando LOUO/LOSO) y aplicaciones en contextos de vida libre. Para facilitar la lectura, la tabla se presenta en orientaci├│n horizontal (landscape).

% Tabla completa en landscape (17 artículos)
\begin{landscape}
\footnotesize
\begin{longtable}{@{}p{1.8cm}p{2.5cm}p{1.8cm}p{1.5cm}p{0.9cm}p{1.3cm}p{2.0cm}p{1.8cm}@{}}
\caption{Análisis comparativo exhaustivo de investigaciones sobre clasificación de sedentarismo y actividad física con wearables}
\label{tab:comparativa_investigaciones_completa} \\
\toprule
\textbf{Autor (A├▒o)} & 
\textbf{Objetivo} & 
\textbf{T├®cnica} & 
\textbf{Datos/ Wearable} & 
\textbf{N} & 
\textbf{Validaci├│n} & 
\textbf{M├®trica Principal} & 
\textbf{Resultado} \\
\midrule
\endfirsthead

\multicolumn{8}{c}{\textit{(Continuaci├│n de Tabla \ref{tab:comparativa_investigaciones_completa})}} \\
\toprule
\textbf{Autor (A├▒o)} & 
\textbf{Objetivo} & 
\textbf{T├®cnica} & 
\textbf{Datos/ Wearable} & 
\textbf{N} & 
\textbf{Validaci├│n} & 
\textbf{M├®trica Principal} & 
\textbf{Resultado} \\
\midrule
\endhead

\midrule
\multicolumn{8}{r}{\textit{Continúa en la siguiente página...}} \\
\endfoot

\bottomrule
\endlastfoot

% ========== FUZZY LOGIC ==========

Capitoli et al. (2025) \cite{Capitoli2025FuzzyXAI} & 
Diagnóstico clínico interpretable con lógica difusa probabilística & 
Fuzzy Logic (interpretable) & 
Datos clínicos hospitalarios & 
Clínico & 
Validación clínica & 
Interpretabilidad 100\% & 
Clínicos controlan proceso diagnóstico. Fuzzy = XAI \\

\midrule

Marashi-Hosseini et al. (2023) \cite{MarashiHosseini2023Dietary} & 
Sistema decisi├│n diet├®tica para pacientes m├║ltiples condiciones cr├│nicas & 
Fuzzy Logic Mamdani (1144 reglas) & 
Datos diet├®ticos + cl├¡nicos & 
100 & 
Validaci├│n vs nutricionistas & 
Precisi├│n & 
97\% precisi├│n. Sistema experto fuzzy riguroso \\

\midrule

Gonçalves et al. (2021) \cite{Goncalves2021} & 
Reconocimiento actividad humana con clustering + fuzzy & 
K-Means ($k=2$) $\rightarrow$ FIS Mamdani & 
Sensores inerciales (IMU) & 
--- & 
Validaci├│n estabilidad & 
Estabilidad clasificaci├│n & 
ÚNICO precedente clustering $\rightarrow$ fuzzy supervisado \\

\midrule

% ========== DEEP LEARNING ==========

Khan et al. (2024) \cite{Khan2024Wearable} & 
Reconocimiento actividad física con sensores inerciales + deep learning & 
Deep Learning (CNNs + LSTMs) & 
Sensores inerciales wearable & 
--- & 
K-fold CV & 
Accuracy & 
Alta precisi├│n. Requiere datos masivos + no interpretable \\

\midrule

Bassani et al. (2025) \cite{Bassani2025DLHAR} & 
Deep Learning para reconocimiento actividad física en tareas manuales & 
Deep Learning (varios modelos) & 
Aceler├│metros & 
--- & 
70/30 vs LOSO & 
Accuracy & 
\textbf{70/30: 95.7\%} $\rightarrow$ \textbf{LOSO: 90.3\%}. Generalizaci├│n robusta requiere LOSO \\

\midrule

% ========== MACHINE LEARNING CLÁSICO ==========

Fuller et al. (2021) \cite{Fuller2021Predicting} & 
Clasificaci├│n lying, sitting, walking, running con Apple Watch + Fitbit & 
Random Forest & 
Apple Watch + Fitbit (aceler├│metro) & 
33 & 
LOOCV & 
Accuracy & 
Accuracy $>$90\% para posturas. BYOD validado \\

\midrule

Mathew et al. (2024) \cite{Mathew2024LOSO} & 
Medición uso funcional mano en parálisis cerebral unilateral con ML & 
Machine Learning (varios modelos) & 
Aceler├│metros (mu├▒eca) & 
22 & 
Personalizado vs LOSO & 
F1-Score & 
\textbf{Personalizado: F1=0.896} $\rightarrow$ \textbf{LOSO: F1=0.584}. Evidencia sobreajuste $N$ peque├▒o \\

\midrule

Rehman et al. (2024) \cite{Rehman2024LOSO} & 
Estudio comparativo validación + extracción características en HAR wearables & 
Random Forest & 
Sensores inerciales wearables & 
--- & 
LOSO vs k-fold & 
Accuracy & 
\textbf{LOSO: 76\%} vs \textbf{k-fold: 89\%}. Data leakage justifica LOUO \\

\midrule

% ========== MÉTODOS ESTADÍSTICOS ==========

Salim et al. (2024) \cite{Salim2024STEPHEN} & 
Detecci├│n tiempo sedentario + bouts con Hidden Semi-Markov Model & 
Hidden Semi-Markov Model & 
Wearables consumo (mu├▒eca) & 
--- & 
Validaci├│n vs verdad operativa & 
Sensibilidad/ Especificidad & 
Detecci├│n precisa sedentary bouts. M├®todo estad├¡stico avanzado \\

\midrule

Migueles et al. (2022) \cite{Migueles2022GRANADA} & 
Consenso análisis acelerómetro para comportamientos físicos en epidemiología & 
Consenso metodol├│gico & 
Aceler├│metros (varios) & 
--- & 
Revisión sistemática & 
Guías metodológicas & 
Estándar internacional para análisis PA/sedentarismo \\

\midrule

% ========== VALIDACIÓN WEARABLES ==========

Lyons et al. (2024) \cite{Lyons2024StandHour} & 
Definici├│n umbrales actividad que activan ``Stand Hour'' en Apple Watch & 
Estudio observacional & 
Apple Watch (energía activa + pasos) & 
--- & 
Análisis umbral & 
Umbrales Stand Hour & 
Deconstruye proxy sedentarismo comercial. Umbrales caja negra identificados \\

\midrule

O'Grady et al. (2024) \cite{OGrady2024AppleWatch} & 
Validaci├│n HRV + RHR en Apple Watch Series 9 y Ultra 2 & 
Estudio validaci├│n & 
Apple Watch vs Polar H10 (gold standard) & 
39 & 
Validaci├│n concurrente & 
MAPE & 
\textbf{HRV MAPE=28.88\%}, \textbf{RHR MAPE=5.91\%}. Limitaciones sensor PPG \\

\midrule

Ji et al. (2023) \cite{Ji2023Scratch} & 
Evaluación scratch nocturno con actigrafía en pacientes dermatitis atópica & 
Machine Learning & 
Actigrafía (ActiGraph) & 
--- & 
\textbf{LOSO evaluation} & 
Sensibilidad/ Especificidad & 
\textbf{Nature npj Digital Medicine}. Precedente metodológico clínico LOSO \\

\midrule

% ========== XAI + INTERPRETABILIDAD ==========

Bienefeld et al. (2023) \cite{Bienefeld2023XAI} & 
Resolver dilema XAI: brecha necesidades clínicos vs objetivos desarrolladores & 
Estudio cualitativo & 
Entrevistas + cuestionarios & 
112 & 
Análisis cualitativo & 
Temas identificados & 
Clínicos requieren transparencia global. Desarrolladores enfocan local \\

\midrule

Deng et al. (2023) \cite{Deng2023LharJBHI} & 
HAR ligero en wearables con knowledge distillation & 
Deep Learning (ligero) & 
Sensores wearables & 
--- & 
Validaci├│n experimental & 
Accuracy + eficiencia & 
\textbf{IEEE JBHI} (revista objetivo). HAR eficiente wearables \\

\midrule

% ========== FISIOLOGÍA + HRV ==========

Godkin et al. (2025) \cite{Godkin2025Context} & 
Medir RHR durante vida diaria: impacto contexto conductual + criterios metodol├│gicos & 
Estudio observacional & 
Wearables (RHR + aceler├│metro) & 
--- & 
Análisis contextual & 
RHR por contexto & 
\textbf{RHR sedentario $\neq$ RHR sue├▒o}. Contextos auton├│micos diferentes. Clasificaci├│n fisiol├│gica vs postural \\

\midrule

Marino et al. (2024) \cite{Marino2024ARIC} & 
Asociaciones PA + HRV con funci├│n cognitiva + demencia (ARIC Neurocognitive Study) & 
Análisis cohorte longitudinal & 
ECG ambulatorio 2 semanas + aceler├│metro & 
961 & 
Análisis regresión & 
Odds Ratio (OR) & 
Mayor PA + HRV $\rightarrow$ mejor cognici├│n. Relaci├│n no lineal PA-HRV-salud cerebral \\

\midrule

% ========== INVESTIGACIÓN ACTUAL (ÚLTIMA FILA) ==========

\textbf{Martínez-Corral (2025)} \newline \textbf{[Investigación Actual]} & 
\textbf{Clasificaci├│n comportamiento sedentario mediante l├│gica difusa + datos biom├®tricos Apple Watch (HRV, aceler├│metro)} & 
\textbf{Clustering K-Means ($k=2$) $\rightarrow$ Fuzzy Mamdani (4 variables)} & 
\textbf{Apple Watch (PPG + IMU) vida libre} & 
\textbf{10 (5F/5M)} & 
\textbf{LOUO (10-fold)} & 
\textbf{F1-Score} & 
\textbf{F1-global=0.840, F1-LOUO=0.780$\pm$0.167. Interpretabilidad 100\%. BYOD validado. Único estudio clustering$\rightarrow$fuzzy + LOUO + HRV sedentarismo} \\

\end{longtable}
\end{landscape}

% ============================================================================
% TABLA DE NOMENCLATURA (ATLAS)
% ============================================================================

\section{Nomenclatura y Simbología Matemática}
\label{sec:nomenclatura}

Esta sección presenta la nomenclatura matemática utilizada en el sistema de inferencia difusa y análisis estadístico.

\subsection{Símbolos Matemáticos Generales}

% Tabla multi-página con longtable (permite continuación automática)
\begin{longtable}{@{}p{3cm}p{10cm}@{}}
\caption{Símbolos y Notación Matemática del Sistema Difuso}
\label{tab:nomenclatura} \\
\toprule
\textbf{Símbolo} & \textbf{Descripción} \\
\midrule
\endfirsthead

\multicolumn{2}{c}{\textit{(Continuaci├│n de Tabla \ref{tab:nomenclatura})}} \\
\toprule
\textbf{Símbolo} & \textbf{Descripción} \\
\midrule
\endhead

\midrule
\multicolumn{2}{r}{\textit{Continúa en la siguiente página...}} \\
\endfoot

\bottomrule
\endlastfoot
\multicolumn{2}{@{}l}{\textit{\textbf{Conjuntos y Funciones de Membresía}}} \\
\(\tilde{A}\) & Conjunto difuso \\
\(\mu_{\tilde{A}}(x)\) & Función de membresía del conjunto difuso \(\tilde{A}\) \\
\(X\) & Universo de discurso (conjunto clásico de valores posibles) \\
\(X_i\) & Variable de entrada \(i\) (\(i \in \{1,2,3,4\}\)) \\
\(Y\) & Variable de salida (Índice de Sedentarismo) \\
\(T_i\) & Conjunto de t├®rminos ling├╝├¡sticos de variable \(i\) \\
\(\mathbf{P}_i\) & Matriz de percentiles \(3 \times 3\) para variable \(X_i\) \\
\(p_{k,i}\) & Percentil \(k\) de la variable \(X_i\) \\
\(\mathbf{M}^{(j)}\) & Matriz de membresías \(4 \times 3\) para observación \(j\) \\
\\
\multicolumn{2}{@{}l}{\textit{\textbf{Variables Fisiol├│gicas}}} \\
\(\text{Act}_{\text{rel}}\) & Actividad Relativa (min movimiento / hrs monitoreadas) \\
\(\text{Sup}_{\text{cal}}\) & Superávit Calórico Basal (\% del TMB) \\
\(\text{HRV-SDNN}\) & Variabilidad Cardíaca (desviación estándar intervalos RR) \\
\(\Delta_{\text{FC}}\) & Delta Cardíaco (FC caminata - FC reposo) \\
\(\text{TMB}\) & Tasa Metabólica Basal (kcal/día, ec. Mifflin-St Jeor) \\
\(\text{GCA}\) & Gasto Calórico Activo (kcal/día) \\
\(\text{FC}\) & Frecuencia Cardíaca (latidos por minuto) \\
\(\text{IQR}\) & Rango Intercuartílico (Q3 - Q1) \\
\\
\multicolumn{2}{@{}l}{\textit{\textbf{Reglas de Inferencia y Activaci├│n}}} \\
\(R_r\) & Regla de inferencia \(r\) (\(r \in \{1,2,3,4,5\}\)) \\
\(\mathcal{R}\) & Base de reglas (conjunto de 5 reglas Mamdani) \\
\(w_r^{(j)}\) & Activaci├│n (firing strength) de regla \(R_r\) para observaci├│n \(j\) \\
\(\mathbf{w}^{(j)}\) & Vector de activaciones \(\in \mathbb{R}^5\) \\
\(T(a, b)\) & T-norm (operador AND fuzzy): \(T(a,b) = \min(a,b)\) \\
\(y_r\) & Output crisp de la regla \(R_r\) \\
\(\mathbf{y}_{\text{salidas}}\) & Vector de salidas de reglas \([1.0, 0.0, 0.9, 0.5, 0.7]^T\) \\
\\
\multicolumn{2}{@{}l}{\textit{\textbf{Defuzzificaci├│n y Clasificaci├│n}}} \\
\(y^{(j)}\) & Score de sedentarismo defuzzificado para semana \(j\) \\
\(\hat{y}^{(j)}\) & Clasificaci├│n binaria (0=Bajo, 1=Alto sedentarismo) \\
\(\tau\) & Umbral de decisi├│n para binarizaci├│n \\
\(\|\mathbf{w}\|_1\) & Norma L1 del vector (suma de componentes) \\
\\
\multicolumn{2}{@{}l}{\textit{\textbf{Validaci├│n LOOU}}} \\
\(\mathcal{D}\) & Dataset completo (1,337 semanas, 10 usuarios) \\
\(\mathcal{D}_{\text{train}}^{(i)}\) & Conjunto de entrenamiento en fold \(i\) (9 usuarios) \\
\(\mathcal{D}_{\text{test}}^{(i)}\) & Conjunto de prueba en fold \(i\) (1 usuario) \\
\(n_i\) & N├║mero de semanas del usuario \(i\) \\
\(F1^{(i)}\) & F1-Score en fold \(i\) (usuario \(i\) como test) \\
\(\overline{F1}_{\text{LOOU}}\) & F1-Score promedio LOOU (10 folds) \\
\(\sigma_{F1}\) & Desviación estándar de F1-Score entre folds \\
\(CV\) & Coeficiente de variaci├│n (\%) \\
\\
\multicolumn{2}{@{}l}{\textit{\textbf{Clustering (Verdad Operativa)}}} \\
\(K\) & N├║mero de clusters (K=2) \\
\(\mathbf{C}_k\) & Centroide del cluster \(k\) (\(k \in \{0,1\}\)) \\
\(c^{(j)}\) & Cluster asignado a semana \(j\) (verdad operativa) \\
\(S\) & Silhouette Score (m├®trica de calidad clustering) \\
\\
\multicolumn{2}{@{}l}{\textit{\textbf{M├®tricas de Evaluaci├│n}}} \\
\(\text{TP}\) & True Positives (verdaderos positivos) \\
\(\text{TN}\) & True Negatives (verdaderos negativos) \\
\(\text{FP}\) & False Positives (falsos positivos) \\
\(\text{FN}\) & False Negatives (falsos negativos) \\
\(\text{Precision}\) & Precisi├│n: \(\text{TP}/(\text{TP}+\text{FP})\) \\
\(\text{Recall}\) & Sensibilidad: \(\text{TP}/(\text{TP}+\text{FN})\) \\
\(F1\) & F1-Score: \(2 \cdot \text{Prec} \cdot \text{Rec} / (\text{Prec} + \text{Rec})\) \\
\(\text{MCC}\) & Matthews Correlation Coefficient \\
\(\text{Acc}\) & Accuracy (exactitud): \((\text{TP}+\text{TN})/N\) \\
\\
\multicolumn{2}{@{}l}{\textit{\textbf{Operadores y Notación Matemática}}} \\
\(\min(a, b)\) & Mínimo de \(a\) y \(b\) \\
\(\max(a, b)\) & Máximo de \(a\) y \(b\) \\
\(\mathbb{E}[\cdot]\) & Valor esperado (esperanza matemática) \\
\(\mathbb{1}[\cdot]\) & Funci├│n indicadora (1 si verdadero, 0 si falso) \\
\(\|\mathbf{v}\|_2\) & Norma euclidiana (L2) del vector \(\mathbf{v}\) \\
\(\|\mathbf{v}\|_1\) & Norma L1 (suma de valores absolutos) \\
\(\mathbb{R}\) & Conjunto de n├║meros reales \\
\(\mathbb{R}^+\) & Conjunto de n├║meros reales positivos \\
\(\mathbb{R}^{n}\) & Espacio euclidiano n-dimensional \\
\([a, b]\) & Intervalo cerrado de \(a\) a \(b\) \\
\(\arg\min_x f(x)\) & Argumento que minimiza la funci├│n \(f(x)\) \\
\(\forall\) & Para todo (cuantificador universal) \\
\(\exists\) & Existe (cuantificador existencial) \\
\(\in\) & Pertenece a (relaci├│n de pertenencia) \\
\(\subset\) & Subconjunto de \\
\(\setminus\) & Diferencia de conjuntos \\
\(\land\) & Conjunci├│n l├│gica (AND) \\
\(\Rightarrow\) & Implicaci├│n l├│gica (THEN) \\
\end{longtable}

% ============================================================================
% FIN DE NOMENCLATURA (ATLAS)
% ============================================================================

% Aquí se incluirán otros anexos de la tesis
% Por ejemplo: consentimiento informado, instrumentos, etc.



